% =============================================================================
% Seção 2 — Trabalhos Relacionados
% =============================================================================

\section{Trabalhos Relacionados}
\label{sec:relacionados}

\subsection{Sistemas de CAA para Doenças Neurodegenerativas}

A comunicação aumentativa e alternativa (CAA) para pessoas com ELA tem sido
estudada extensivamente desde a década de 1990~\cite{beukelman2013aac}. Soluções
comerciais dominam a prática clínica: dispositivos de rastreamento ocular de ponta,
como o Tobii Dynavox PCEye e o Eyegaze Edge, oferecem seleção por olhar com precisão
superior a 98\,\% em condições laboratoriais, mas apresentam custos de hardware acima
de USD~3.000 e requerem recalibração periódica por profissionais treinados~\cite{beukelman2013aac}.
Esses sistemas pressupõem controle oculomotor residual suficiente, que se deteriora
à medida que a ELA progride em direção à oftalmoplegia. Pacotes de software de CAA
como o Grid~3 dependem de conectividade Bluetooth ou Wi-Fi estável para sincronizar
conjuntos de símbolos e perfis de usuário, criando uma barreira para usuários
domiciliares sem banda larga confiável.

Uma pesquisa de Clarke et al.~\cite{clarke2021aac} identificou que a provisão de CAA
para pessoas com ELA é caracterizada por encaminhamento tardio, transições rápidas de
dispositivos à medida que a função motora declina, e uma lacuna persistente em soluções
de baixo custo e adaptáveis adequadas para o ambiente domiciliar. As interfaces
cérebro-computador (BCI) via eletroencefalografia oferecem um caminho alternativo que
contorna o sistema motor por completo, mas os headsets de BCI não invasivos atuais
atingem taxas típicas de transferência de informação de 5--25 bits por minuto em
condições ótimas, requerem operadores treinados e são inadequados para o contexto de
cuidado domiciliar~\cite{zander2011passive}. BCIs corticais implantadas atingem taxas
maiores, mas apresentam risco cirúrgico. O NeuroAssist AI visa o estágio intermediário
— após a perda da função dos membros e da fala, mas enquanto persiste algum grau de
movimento facial voluntário — para o qual não existe atualmente nenhuma solução de
baixo custo e independente de rede na literatura.

\subsection{Detecção de Gesto Facial via Visão Computacional}

Três gerações de análise facial baseada em câmera são relevantes para este trabalho.

\emph{Primeira geração — classificadores em cascata.} Viola e Jones~\cite{violajones2004}
introduziram a detecção de faces em tempo real usando cascatas de características Haar,
permitindo a implantação em hardware de consumo. Esses detectores localizam o rosto,
mas não fornecem coordenadas de marcos sub-faciais; a detecção de piscadas requer
heurísticas adicionais aplicadas à região do olho detectada e é altamente sensível a
variações de iluminação.

\emph{Segunda geração — regressão de marcos faciais.} A biblioteca Dlib de King~\cite{king2009dlib}
popularizou a detecção de 68 marcos faciais por meio de um conjunto de árvores de
regressão, permitindo o cálculo confiável do EAR~\cite{soukupova2016} sem aprendizado
profundo. Soukupová e Čech~\cite{soukupova2016} formalizaram o EAR como um escalar
único para detecção de piscadas, demonstrando precisão de detecção em nível de quadro
acima de 99\,\% em conjuntos de dados controlados. A principal limitação dessa geração
é a ausência de um runtime compatível com WebAssembly, tornando a implantação no
navegador do lado do cliente impraticável.

\emph{Terceira geração — modelos de marcos profundos no dispositivo.} O MediaPipe
FaceLandmarker~\cite{lugaresi2019mediapipe} fornece 468 marcos faciais normalizados e
52 coeficientes de blendshape a aproximadamente 30\,fps em um contexto de navegador
WebAssembly usando apenas um delegate de CPU. O NeuroAssist AI se baseia nessa
geração, retendo o EAR para a máquina de estados de dupla piscada (onde sua
simplicidade auxilia a interpretabilidade) enquanto usa pontuações de blendshape para
gestos de piscada de um olho e sobrancelha.

\subsection{Reconhecimento de Gesto Deliberado e Supressão de Falsos Positivos}

Um desafio fundamental em CAA acionada por gesto é distinguir ativações intencionais
da atividade facial involuntária que ocorre continuamente durante o comportamento normal.
A taxa de piscadas espontâneas em adultos saudáveis é de 15--20 piscadas por
minuto~\cite{bentivoglio1997}; em pacientes com ELA, as taxas podem ser menores devido
à perda de neurônios motores, mas permanecem diferentes de zero. Qualquer abordagem
baseada em limiar aplicada no nível de quadro único produz taxas de falsos positivos
da ordem de vários por minuto em condições naturais, tornando o monitoramento sem
supervisão inviável.

As duas principais estratégias na literatura para suprimir falsos positivos em CAA
facial são \emph{tempo de permanência} (\emph{dwell time}) e \emph{modelagem de
intencionalidade}. A seleção por tempo de permanência — amplamente usada em CAA
baseada em olhar — exige que o usuário fixe uma região-alvo por uma duração fixa
(tipicamente 0,5--1\,s)~\cite{majaranta2002eyetyping}. Essa abordagem é inaplicável
à entrada baseada em piscada, pois uma piscada por definição fecha o olho e interrompe
a fixação. A modelagem de intencionalidade via máquinas de estados foi proposta na
literatura mais ampla de interação humano-computador~\cite{zander2011passive}, mas sua
aplicação à CAA baseada em piscada para ELA não foi sistematicamente descrita. A
máquina de seis estados introduzida neste artigo (Seção~\ref{sec:maquinadeestados})
aplica restrições de intencionalidade específicas ao gesto de dupla piscada: duração
mínima de fechamento por piscada, abertura confirmada entre piscadas, e uma janela
temporal — combinação não encontrada, ao conhecimento dos autores, na literatura
anterior de CAA.

\subsection{Computação de Borda e Privacidade para Aplicações de Saúde}

O processamento de dados de sensores relacionados à saúde no dispositivo, em vez de
na nuvem, tem recebido atenção crescente após a aplicação do GDPR em 2018 e da LGPD
em 2020~\cite{edpb2020privacy}. O vídeo facial é classificado explicitamente como dado
biométrico sob ambos os marcos regulatórios, atraindo as maiores obrigações regulatórias
de processamento. A execução local via WebAssembly fornece uma garantia prática de
privacidade: o modelo de execução com memória isolada da especificação impede que um
módulo carregado acesse qualquer estado do navegador fora de sua memória linear, e
nenhuma chamada de rede é necessária para executar a inferência do marcador de
landmark~\cite{haas2017wasm}. A latência de implantação de modelos WebAssembly para
análise facial foi relatada em 20--40\,ms por quadro em hardware de consumo sem
aceleração por GPU~\cite{lugaresi2019mediapipe}, consistente com a operação a 30\,fps
observada em nosso PoC.

\subsection{Arquitetura de Software para Sistemas de Saúde}

A arquitetura Hexagonal (Portas e Adaptadores)~\cite{cockburn2005hexagonal} foi
proposta como um padrão estrutural para sistemas de informação em saúde onde o modelo
de domínio deve permanecer isolado de mudanças tecnológicas — um requisito recorrente
em software de saúde regulamentado onde requisitos de auditoria, esquemas de dados e
padrões de integração evoluem de forma independente~\cite{richardson2018microservices}.
A decomposição em microsserviços de plataformas de CAA não foi relatada na literatura,
ao conhecimento dos autores; a maioria dos sistemas existentes são aplicações
monolíticas de desktop ou firmware. A arquitetura do NeuroAssist AI demonstra que
decompor um pipeline de gesto-para-alarme em serviços independentemente implantáveis
é viável em escala de PoC, e que a aplicação de limites em tempo de compilação via
ArchUnit fornece um mecanismo prático para manter a disciplina hexagonal.

\subsection{Comunicação e Segurança Durante o Uso de VNI na ELA}

A adesão à ventilação não invasiva na ELA está consistentemente abaixo das metas
clínicas, com desconforto relatado pelo paciente e assincronia dispositivo-paciente
identificados como principais fatores contribuintes~\cite{berlowitz2021niv, yang2026pva}.
A polissonografia domiciliar em 23 pacientes com ELA utilizando VNI noturna documentou
um índice de assincronia médio de 69 eventos por hora (intervalo 15--146) — muito
acima dos limiares associados à perturbação clinicamente significativa do
sono~\cite{atkeson2011pva}. O ensaio clínico randomizado 3TLA~\cite{berlowitz2025tla}
investiga atualmente se a titulação da VNI guiada por polissonografia melhora as taxas
de uso; seu desfecho primário — uso de VNI superior a quatro horas por dia — reflete
o desafio de manter a adesão em ambientes domiciliares. Uma dimensão complementar e
pouco explorada é a capacidade do paciente de sinalizar desconforto ou angústia aos
cuidadores durante o uso de VNI sem supervisão, particularmente à noite. Ao conhecimento
dos autores, nenhum sistema publicado fornece um mecanismo independente de rede e de
baixo custo para esse fim; o NeuroAssist AI foi projetado para preencher essa lacuna.
